\documentclass[11pt,a4paper]{article}
\usepackage[margin=27mm]{geometry}
\usepackage{amsmath,amssymb,amsthm,booktabs,longtable,hyperref}
\hypersetup{hidelinks}
\setlength{\emergencystretch}{2em}
\newtheorem{theorem}{Theorem}
\newtheorem{lemma}[theorem]{Lemma}
\newcommand{\im}{\operatorname{im}}
\newcommand{\Aut}{\operatorname{Aut}}
\newcommand{\Soc}{\operatorname{Soc}}
\title{Ternary codes from the Witt design and the plane of order four}
\author{Private application manuscript: independent review pending}
\date{}
\begin{document}\maketitle
\begin{abstract}
We derive the ternary graph codes studied by Leemans and Rodrigues from
one incidence calculation and the common criteria for finite module and
automorphism certificates. The complete weight records and actual
permutation actions determine all four principal codes, their specified
submodules and full coordinate automorphism groups. We also give the
prime-rank table, the socle and auxiliary-module calculations, and the
precise qualifications needed for several ancillary source statements.
Classifications of other irreducible representations remain explicitly
retained inputs, distinct from the irreducibility certificates for the
modules constructed here.
\end{abstract}

The source is the complete eleven-page 2016 edition of
Leemans--Rodrigues~\cite{LR}. All four clauses of its main Theorem~1 are
addressed, together with Lemmas~1--2, Propositions~1--5, Corollary~1,
Tables~1--2 and the specified modules and orbits in Remarks~1--3.
We give the corrected meaning of the ancillary assertions that cannot
hold literally. The common paper~\cite{U} supplies full proofs of the
projection, incidence, matrix-algebra, socle, refinement and weight
identities used below. Its whole cost belongs to the topic; it is not a
free import from another application. Every finite certificate is bound
to the displayed construction and its actual basis. The programme's
numerical comparisons and source admissions are separate from this draft.

\section{The four graph codes}

Throughout, automorphisms of a code mean permutations of its coordinates,
not arbitrary monomial transformations. Unless a prime is displayed, the
field is $F=\mathbb F_3$. The constant vector is denoted by $j$, with its
length determined by context. A direct sum of submodules below is an
internal sum in the same coordinate space, rather than a concatenation of
shorter codes. For a symmetric adjacency matrix $A_n$ of order $n$, put
\[
 C_n=\im A_n,\qquad
 D_n=\langle (A_n)_i-(A_n)_1:1\leq i\leq n\rangle.
\]
For length $120$ we also write $S=D_{120}$, to agree with the source.
The code of the complementary graph is denoted by $\overline C_n$.
This notation distinguishes a graph complement from an orthogonal dual.

The main results of Leemans and Rodrigues~\cite{LR} concern four specified
strongly regular graphs. We use their graph constructions and the
identifications with the relevant McLaughlin subgraphs, as attributed in
that paper to Brouwer--van Eijl and the preceding graph literature.
The actual graphs are fixed here, so assertions about a particular code
are not inferred merely from four strongly regular parameters.

For $n=105$, use the incident point--line flags of $\operatorname{PG}(2,4)$.
Join distinct flags $(p,\ell),(q,m)$ if $p=q$, or $\ell=m$, or neither
cross incidence $p\in m$, $q\in\ell$ holds. The finite field is
$\mathbb F_4=\mathbb F_2[z]/(z^2+z+1)$. All $105$ vertices and their
adjacencies are constructed directly from this field. The printed
adjacency sentence on page~242 has a defect: a conjunctive reading gives
valency eight, while joining all its clauses disjunctively gives forty.
Our identification with the intended graph therefore retains Coolsaet's
uniqueness theorem for the complementary parameters $(105,32,4,12)$
\cite[Theorem~3]{Coolsaet}; the verified complementary parameters alone
would not identify a graph without that input.

For the other three lengths, let $\mathcal D$ be the specified
$S(4,7,23)$ with its $253$ heptads. One complete construction takes the
weight-seven words in the binary span of the twelve shifts of
\[
 1+x+x^5+x^6+x^7+x^9+x^{11}.
\]
The supplied constructor verifies all $4096$ words, their weight
multiplicities, and the unique containment of every four-subset in a
heptad. It also verifies all block intersection sizes. Delete a fixed set
of $t$ points, where $t=0,1,2$, and retain the heptads avoiding it. There
are respectively $253,176,120$ such blocks. Join distinct retained blocks
when their intersection has size three. This is exactly the complement
of the source's intersection-one graph. With $H_t$ the point--block
incidence matrix on the remaining $23-t$ points, the common paper
\cite[Theorem~2.2]{U} gives
\[
\begin{array}{c|c|c|c}
 t&n&H_tH_t^{\mathsf T}\pmod3&C_n\\\hline
 0&253&2I_{23}&\im H_0^{\mathsf T}\\
 1&176&I_{22}+J_{22}&H_1^{\mathsf T}(e^\perp)\\
 2&120&I_{21}&\im H_2^{\mathsf T}.
\end{array}                                                    \tag{1}
\]
Here $e$ is the constant point vector. Each of these three codes is LCD;
$\overline C_n=C_n^\perp$ for $n=120,176$. Also
$C_{253}=Fj\mathbin{\perp\!\oplus}D_{253}$, with
$D_{253}=H_0^{\mathsf T}(e^\perp)$ of dimension $22$ and LCD.
The hypotheses of the common theorem are discharged by the complete
Steiner and intersection checks, not by assuming the ternary ranks.

The verified parameters and integral spectra are
\[
\begin{array}{c|c|c}
n&(n,k,\lambda,\mu)&\text{spectrum of }A_n\\\hline
105&(105,72,51,45)&72^1,9^{20},(-3)^{84}\\
120&(120,77,52,44)&77^1,11^{20},(-3)^{99}\\
176&(176,105,68,54)&105^1,17^{21},(-3)^{154}\\
253&(253,140,87,65)&140^1,25^{22},(-3)^{230}.
\end{array}                                                     \tag{2}
\]
The finite certificate checks every vertex degree and every pair's
common-neighbour count. The spectra then follow from
$A^2=(\lambda-\mu)A+(k-\mu)I+\mu J$, the constant eigenvector,
and the trace and dimension identities. Connectivity follows from
$\mu>0$, so the constant eigenvalue has multiplicity one.

\subsection{The retained orbital construction}
The source's Result~1 attributes a general construction to Key and Moori,
including their correction. Its argument is short enough to make the
exact scope explicit. Let $G$ act primitively on $\Omega$, let $\alpha\in
\Omega$, and let $\Delta\ne\{\alpha\}$ be a nonempty orbit of $G_\alpha$.
The set stabilizer $G_\Delta$ contains $G_\alpha$. Primitivity makes the
latter maximal; since $\Delta$ is a proper nonempty subset,
$G_\Delta\ne G$. Hence $G_\Delta=G_\alpha$, and the translates of $\Delta$
are $|\Omega|$ distinct blocks. Transitivity and counting incidences give
constant block size and replication $|\Delta|$, producing the symmetric
$1$-design. If the orbital is self-paired, these blocks are precisely the
neighbour sets of the associated undirected graph. The nonempty graph's
components form a $G$-invariant partition, so it is connected. Its
symmetric incidence matrix identifies the design with its dual, and both
point and block actions are primitive. This preserves the background
result without using it in place of any of the explicit constructions above.

\section{Ranks and the complete finite weight records}

The two prime-rank lemmas in the source, and its surrounding assertion
about other characteristics, can be handled together. For the fixed
matrices the full answer is
\[
\begin{array}{c|l|c}
n&\text{exceptional primes and ranks}&\text{rank at every other prime}\\\hline
105&2:104,\quad3:20&105\\
120&3:21,\quad7:119,\quad11:99&120\\
176&3:21,\quad5:175,\quad7:175,\quad17:155&176\\
253&2:252,\quad3:23,\quad5:230,\quad7:252&253.
\end{array}                                                     \tag{3}
\]
Indeed (2) determines the integer determinant as the product of the
listed eigenvalues with their multiplicities. Outside its prime divisors
reduction is invertible. For every prime divisor, the complete supplied
elimination checks the rank of the actual matrix. Thus the finite table
covers every exceptional characteristic; no sample of small primes is
being extrapolated. For the ternary entries, (1) also gives a written
incidence proof for $n=120,176,253$, while the $105$ entry is certified
by a complete row reduction.

This proves Lemmas~1 and~2, including the full ranks at $p=11,23$ in
Lemma~2. The paragraph preceding Lemma~1 says that the only nonternary
exception to full rank or codimension one occurs for $(n,p)=(253,5)$.
The additional entries $(120,11)$ and $(176,17)$ in (3) correct that
sentence. They do not change any of the ternary principal codes.

The distances used below come from the previously completed full
coefficient enumerations in the exact row-reduced bases, together with
the two separately completed enumerations for $K$ and $N$ at length $120$.
These are mathematical input certificates. For a basis of dimension $k$,
the enumerator visits each of the $3^k$ coefficient vectors once, and
basis independence makes their images all distinct codewords. The
complete distributions are retained, not merely their minimum entries.
The current checker compares the reconstructed bases with those exact
bases. No large enumeration is repeated during assembly.

The common MacWilliams theorem \cite[Theorem~7.1]{U} then gives every dual
coefficient over the integers. In particular the resulting parameters
and minimum multiplicities are
\[
\begin{array}{c|r|r|r|r|r|r}
\text{code}&n&k&d&A_d&d^\perp&A_{d^\perp}^\perp\\\hline
C_{105}&105&20&33&770&6&15680\\
D_{105}&105&19&36&3360&5&84\\
C_{120}&120&21&30&112&6&33600\\
S&120&20&30&112&6&43680\\
K&120&15&48&1260&4&2100\\
N&120&16&48&1260&4&1260\\
C_{176}&176&21&56&924&6&184800\\
C_{253}&253&23&77&46&6&850080\\
D_{253}&253&22&97&3542&6&850080.
\end{array}                                                     \tag{4}
\]
The dual dimension is $n-k$ in each row. Nonnegative integer coefficients
and the dual total are checked, but those consistency checks alone do
not certify an arbitrary proposed input distribution. The complete
coefficient enumeration, its matrix basis, and its method are all retained.
The two $105$ distributions agree entry by entry with source Tables~1--2;
those full distributions are reproduced in the data table accompanying
this application. The two further $20$-dimensional codes of Remark~2 use
a smaller, justified search, described below.

\section{The codes on 105 flags}

The graph is the one fixed in Section~1. Here write $C=C_{105}$ and
$D=D_{105}$.

Let $A$ be the adjacency matrix over $F=\mathbb F_3$, and write
\[
 C=\langle\text{rows of }A\rangle,
 \qquad D=\langle A_i-A_1:1\leq i\leq105\rangle.
\]
The parameter identity gives $A^2=0$, so $C$ and $D$ are self-orthogonal.
Exact row reduction gives $\dim C=20$, $\dim D=19$, and
\[
 C=D\oplus Fj,                                      \tag{5}
\]
where $j$ is the constant vector. Both bases agree byte for byte with the
separately completed enumeration. The two complete distributions give
\[
\begin{array}{c|rrrr}
 &\dim&d&\text{minimum words}&(\dim C^\perp,d^\perp,\text{minimum words})\\\hline
 C&20&33&770&(85,6,15680)\\
 D&19&36&3360&(86,5,84).
\end{array}
\]
The complete dual distributions are checked again by the classical
MacWilliams identity. No $3^{20}$ or $3^{19}$ enumeration is repeated.
The retained distributions agree entry by entry with the two printed
tables.

\subsection{The full group and irreducibility}
The literal coordinate permutations in the certificate preserve every
entry of $A$. Their closure has order $241920$, and their action has
rank four. They generate a group $G$ with a normal subgroup
$N\cong\operatorname{PSL}_3(4)$ of order $20160$ and an explicit
complement $Q\cong D_{12}$ of order twelve. The complement consists of
a diagonal projectivity, field conjugation and point--line duality.
The checker derives their actual flag permutations, verifies
$N\cap Q=1$ and $NQ=G$, and checks the group identifications.
Thus $G\cong L_3(4):D_{12}$, with this action fixed.

The action on the row-reduced $19$-dimensional basis of $D$ is checked
directly. The generator matrices span a matrix algebra of dimension
$19^2=361$: the verifier evaluates independent generator words and
row-reduces their flattened matrices. The common matrix-algebra criterion \cite[Theorem~6.2]{U} therefore proves absolute irreducibility of $D$. Consequently (5) is an internal decomposition into two simple
modules, as required in Proposition~1.

The complete minimum-word set of $C$ is obtained as four disjoint orbits
of literal weight-$33$ words, of lengths $105,105,280,280$.
The accepted count $770$ proves exhaustion, and also recovers Remark~1.
There are $385$ distinct supports. As for the $120$- and $176$-coordinate
codes, form their bipartite coordinate--support incidence graph.
Apply the common refinement bound \cite[Theorem~6.4]{U}. 
Successively individualizing
coordinates $1,2,3,4,7,12,6$ gives cell sizes
\[
 105,8,3,2,4,3,4.
\]
All coordinate colours then differ, so the theorem gives
\[
 |\Aut(C)|\leq105\cdot8\cdot3\cdot2\cdot4\cdot3\cdot4=241920.
\]
Since $G\leq\Aut(\Gamma)\leq\Aut(C)$, equality follows.
Every automorphism of $D$ fixes $j$ and hence preserves $D+Fj=C$.
Conversely $G$ preserves $D$, so
\[
 \Aut(D)=\Aut(C)=\Aut(\Gamma)=L_3(4):D_{12}.          \tag{6}
\]
These are coordinate-permutation groups, not monomial groups.

Both retained distributions have exactly $3360$ words of weight $36$.
Since $D\subset C$, all such words of $C$ are in $D$.
Their span is a nonzero $G$-submodule of the irreducible $D$, and is
therefore $D$. This establishes the source's intrinsic definition of its
codimension-one code, rather than identifying it only by parameters.

\subsection{Uniqueness inside the permutation module}
Put $M=F^{105}$. A supplied basis gives invariant prefixes of dimensions
\[
 0,1,20,39,84,85,104,105.
\]
The factor dimensions are $1,19,19,45,1,19,1$.
All factors are simple. For the factors of dimensions one and nineteen,
this is verified by the full matrix algebra as above. For the $45$-dimensional factor use the rank-one alternative in that same theorem.

Here a recorded polynomial of a recorded sum of generator words gives
$P$ of rank one. The checker evaluates it and produces both bases of
size $45$. Hence its algebra has dimension $2025$, proving simplicity
over $F$ and after extension. This is a certificate of the concrete
factor, not an assumed classification of representations.

For each of the three $19$-dimensional factors $S_i$, all linear equations
\[
 XP_g=(S_i)_gX
\]
are solved, where $P_g$ is the actual $105$-coordinate permutation
matrix. The full solution-space dimensions are $0,1,1$ in the displayed
filtration order. Every basis solution has all rows in $D$, hence every
homomorphism $S_i\to M$ has image in $D$.

Any $19$-dimensional submodule of $M$ is simple. Indeed its composition
factors occur among those of $M$; dimension $45$ is impossible, while
the three available one-dimensional factors alone cannot account for
dimension nineteen. Its sole factor must therefore have dimension
nineteen. Such a simple submodule is isomorphic to one of the $S_i$,
and its inclusion is the image of a corresponding homomorphism into
$M$. The preceding complete equations force its image to be $D$.
Thus $D$ is the unique submodule of dimension nineteen in this
permutation representation, proving the first assertion of Remark~2.

\subsection{The other two twenty-dimensional modules}
The certificate supplies invariant row-reduced bases for $C_1,C_2$,
besides $C_3=C$. Their containments are
\[
\begin{array}{c|cc}
 &D\subseteq C_i&j\in C_i\\\hline
C_1&\text{no}&\text{yes}\\
C_2&\text{yes}&\text{no}\\
C_3&\text{yes}&\text{yes}.
\end{array}
\]
The only one-dimensional submodule of $M$ is $Fj$. To see this directly,
an invariant line has a spanning vector with entries in $\{0,a,-a\}$.
Transitivity makes its zero set empty. A nontrivial scalar action would
interchange the two nonzero values, forcing their multiplicities to be
equal, contrary to the odd length $105$. Thus the line is fixed and
constant. Every $20$-dimensional submodule has factors of dimensions
$19$ and one, by the full composition series above. The containments
and uniqueness just proved give
\[
 \dim\Soc(C_1)=1,\quad \dim\Soc(C_2)=19,\quad
 \dim\Soc(C_3)=20.
\]
The three modules are therefore reducible and pairwise nonisomorphic
as $FG$-modules, not just inequivalent as codes.

For their remaining code parameters, apply the common transitive-distance and minimum-set theorems \cite[Theorems~8.1--8.2]{U} to the actual transitive group and row-reduced information columns.
For $(n,k)=(105,20)$ and candidate distances $30,28$, both the smaller-word
and minimum-word tests require only messages of weight at most five.
There are
\[
 \sum_{i=0}^{5}\binom{20}{i}2^i=583569
\]
such ternary messages, including zero. The complete shared enumerator
checks them once each. Its smallest nonzero output weights are $30,28$,
and literal witnesses attain those values. The averaging theorem
therefore proves those minimum distances.

The unions of a literal witness orbit and its negative have sizes $672$
and $720$. They are invariant sets of minimum words. Exactly $342$ and
$396$ of these words, respectively, have message weight at most five;
these equal the complete restricted enumeration counts at the minimum
weight. Thus every restricted minimum word is in the proposed invariant
set. Averaging an arbitrary minimum word into that restriction proves
that the sets are complete. This recovers
$C_1=[105,20,30]_3$ with $672$ minimum words and
$C_2=[105,20,28]_3$ with $720$ minimum words, as printed in Remark~2.
It neither assumes a complete $3^{20}$ enumeration nor asserts that a
short search is the full weight distribution.

\subsection{The global and local representation claims}
The source's proof of Corollary~1 does not justify its global statement
that nineteen is the smallest dimension of a nontrivial irreducible
$F_3G$-module. Literally that statement is false: $G$ has a normal
subgroup of index two, so its quotient acts nontrivially by the scalars
$1,-1$ on a one-dimensional $F_3$-space. The certificate exhibits the
index-two subgroup by adjoining the diagonal and field automorphisms
to $N$; its order is $120960$. No faithfulness hypothesis is stated in
the printed claim. We prove the precise local replacement: the unique
nonconstant simple submodule of dimension nineteen in this permutation
module is $D$, and no nonconstant simple submodule has smaller dimension.
No classification of all faithful irreducible representations is inferred.

The graph uniqueness input and all finite premises remain those stated above. Source comparisons to historical optimal-code tables are treated in the final section. Discovery used GRAPE and MeatAxe, but the final checks use literal generators, a full invariant filtration, matrix products and complete linear equations.


\section{The graph and difference codes of length 120}

For the residual incidence matrix $H=H_2$, equation (1) gives
$HH^{\mathsf T}=I_{21}$ and $C_{120}=\im H^{\mathsf T}$. Thus the
incidence map is an isometry from the point space to this code. The
constant vector satisfies $H^{\mathsf T}e=j$ and $Hj=40e=e$ over $F$,
so $j\in C_{120}$. The common complement theorem proves
\[
 C_{120}\mathbin{\perp\!\oplus}C_{120}^{\perp}=F^{120},
 \qquad C_{120}^{\perp}=\overline C_{120}.
\]
The $21$-dimensional code and its $99$-dimensional dual have the first
length-$120$ parameters and counts in (4).

The supplementary subcode $S$ in the proof of Proposition~2 is exactly
\[
 S=C_{120}\cap j^\perp=H^{\mathsf T}(e^\perp)
   =\langle(A_{120})_i-(A_{120})_1:1\leq i\leq120\rangle.       \tag{7}
\]
To see the last equality, every row difference has sum zero, while the
rows of $A_{120}$ have the same nonzero sum $77=2$. Their span has
dimension $21$, and adjoining any one row to the difference span gives
that whole space. The difference span therefore has dimension at least
$20$, and containment in the kernel of the nonzero sum functional gives
equality. Under the isometry, that kernel is $e^\perp$.
Since $e^{\mathsf T}e=21=0$, the radical of $e^\perp$ is precisely $Fe$.
Consequently
\[
 \operatorname{Hull}(S)=Fj.                                   \tag{8}
\]
The accepted exact difference-code basis and distribution prove
$S=[120,20,30]_3$ and $S^\perp=[120,100,6]_3$, with respectively
$112$ and $43680$ minimum words.

There is also a direct proof of the source's reflexive-code identity.
On $C_{120}=H^{\mathsf T}F^{21}$ the operator $A_{120}+I$ acts as
$H^{\mathsf T}J_{21}$ on point coordinates; on $C_{120}^\perp=\ker H$
it acts as the identity. Its image is therefore
\[
 \im(A_{120}+I)=Fj+C_{120}^\perp=S^\perp.                     \tag{9}
\]
The two summands have dimensions one and $99$ and meet trivially because
$C_{120}$ is LCD. Thus (9) is an equality of the actual codes, not a match
of parameters.

% The full local group bridge is retained below.
\subsection{Full automorphism groups of the graph and difference codes}

Let $\Gamma$ be the fixed length-$120$ graph, write $A=A_{120}$ and $C=C_{120}$, and let $G$ be the supplied group on its coordinates.

\begin{theorem}\label{L050.120.main}
For these actual coordinate spaces,
\[
 \Aut(C)=\Aut(S)=\Aut(\Gamma)=G,\qquad |G|=80640.
\]
Here $G$ is the group on the literal four coordinate generators of the retained degree-$120$ module certificate. Its previously checked abstract identification is $L_3(4):2^2$. The equality and order in this theorem have the separate elementary certificate below; the abstract identification uses that earlier exact group certificate.
\end{theorem}

\begin{lemma}[Recovering a graph from its code]\label{L050.120.projection}
For this matrix, $E=J-A$ is the orthogonal projector onto $C$. Consequently every coordinate permutation preserving $C$ preserves $A$.
\end{lemma}
\begin{proof}
The strongly regular parameters give $A^2=33I+8A+44J$, hence $A^2=-A-J$ over $\mathbb F_3$. Also $AJ=JA=2J$ and $J^2=0$. Therefore
\[
 E^2=J^2-JA-AJ+A^2=J-A=E,\qquad EA=A.
\]
Every row of $J$ lies in $C$: the sum of the rows of $A$ is $77\mathbf1=2\mathbf1$. Thus $\operatorname{im}E\subseteq C$, and $EA=A$ gives the reverse inclusion (using symmetry to identify row and column spaces). Since $E$ is symmetric and idempotent, its kernel is $C^\perp$ and its image is $C$, so $C\cap C^\perp=0$. This determines the projector uniquely. A coordinate permutation preserving $C$ also preserves the standard dot product and hence $C^\perp$; conjugation must fix $E$. It fixes $J$, so it fixes $A=J-E$.
\end{proof}


\begin{lemma}[Recovering the graph from the difference code]\label{L050.120.support}
The collection of supports of the minimum words of $S$ determines $A$.
\end{lemma}
\begin{proof}
The literal generators of $G$ preserve $A$, hence preserve both $C$ and $S$. Indeed, a permuted row difference is again a difference of two rows. The exact replay verifies a supplied word $w$ lies in the actual $S$ and has weight 30. Its orbit under these generators consists of 112 distinct words. The complete distribution in (4) gives exactly 112 minimum words, so this orbit is all the minimum words of $S$. It is closed under negation and has exactly 56 distinct supports. No orbit size is substituted for an unproved completeness assertion: that assertion uses the retained exhaustive count.

Let $\mathcal B$ be these 56 supports and put
\[
 t_{ij}=|\{B\in\mathcal B:i,j\in B\}|.
\]
The replay constructs every orbit word from $w$, then computes every entry of this integer incidence Gram matrix. It gives $t_{ii}=14$, and, on the 7140 unordered distinct pairs, the complete table is
\begin{center}\begin{tabular}{rrrr}\toprule
$A_{ij}$&$t_{ij}$&Number of pairs&Interpretation\\\midrule
0&4&2520&nonedge\\
1&2&3360&edge\\
1&6&1260&edge\\\bottomrule
\end{tabular}\end{center}
Thus, for $i\ne j$, $A_{ij}=1$ if and only if $t_{ij}\ne4$. Any coordinate permutation preserving $S$ preserves Hamming weight and therefore permutes $\mathcal B$; it preserves every $t_{ij}$ and hence $A$.
\end{proof}

\begin{lemma}[An elementary full graph bound]\label{L050.120.fullness}
The retained four permutations generate all of $\Aut(\Gamma)$.
\end{lemma}
\begin{proof}
First the replay verifies that each is a permutation of all 120 coordinates preserving every adjacency entry. Closing the generated permutations under composition gives exactly 80640 distinct elements; only exact finite permutation operations are used.

For the reverse inequality, apply the common refinement theorem \cite[Theorem~6.4]{U}, using the deterministic choice of the first vertex in the smallest nonsingleton cell (lexicographic ties). 
The following are the actual chosen vertices and sizes of their cells just before individualization; vertices in this table are one-based.
\begin{center}\begin{tabular}{rrrr}\toprule
Step&Chosen vertex&Cell size&Previously fixed vertices\\\midrule
1&1&120&none\\
2&5&42&1\\
3&2&8&1,5\\
4&3&2&1,5,2\\\bottomrule
\end{tabular}\end{center}
After the last refinement all 120 cells are singletons. The common theorem gives $|\Aut(\Gamma)|\le120\cdot42\cdot8\cdot2=80640$.
This matches the subgroup lower bound.
\end{proof}


\begin{proof}[Proof of Theorem~\ref{L050.120.main}]
Every graph automorphism permutes rows of $A$ and therefore preserves $C$ and $S$. Lemmas~\ref{L050.120.projection} and~\ref{L050.120.support} give the reverse inclusions. Lemma~\ref{L050.120.fullness} identifies the common permutation group exactly.
\end{proof}



\section{The socle on 120 Fano subplanes}
Use the fixed $S(4,7,23)$ already constructed and checked by
\texttt{steiner\_construction.py}. Delete two named points. The $120$ heptads
avoiding them form a set $\mathcal H$. The $21$ heptads containing both
points, with those two removed, form the lines of a projective plane on
$21$ points. Its lines have size five; every two remaining points lie on
one line by the Steiner property. Two such lines meet in one point, since
the original heptads meet in one or three points and already share the
deleted pair.

For $H\in\mathcal H$, every two points of $H$ lie on a plane line meeting
$H$ in three points. Indeed the heptad defining that line shares at least
two points with $H$ and hence exactly three. These triples form the Fano
subplane on the seven points of $H$. Let $\mathcal F$ be the $105$ incident
point--line flags of the plane, and define a $105\times120$ matrix $R$ by
\[
 R_{(p,\ell),H}=1\quad\Longleftrightarrow\quad
 p\in H\ \hbox{ and }\ |\ell\cap H|=3.
\]
Thus each column records the $21$ flags of the corresponding Fano subplane.
All following code spaces and matrices are over $\mathbb F_3$.

The common square-zero incidence lemma \cite[Lemma~6.1]{U} will give self-orthogonality after the following integer table is checked.

For the fixed family, the exact integer scalar products of incidence rows
are as follows. This table is checked for every pair of the $105$ flags by
the retained complete construction, not inferred from a sampled orbit.
\[
\begin{array}{l|r}
\text{relation between the two flags}&(RR^{\mathsf T})_{f,g}\\ \hline
f=g&24\\
\text{same point or same line, but distinct}&12\\
\text{distinct points and lines, with one cross incidence}&6\\
\text{distinct points and lines, with neither cross incidence}&3.
\end{array}
\]
Two distinct points and distinct lines cannot have both cross incidences.
Thus the table exhausts all cases and proves $RR^{\mathsf T}=0$ over
$\mathbb F_3$. Each row has sum $24$, so $Rj=0$, and the constant vector on
$120$ coordinates has norm zero. Set
\[
 K=\im(R^{\mathsf T}R),\qquad N=K+\mathbb F_3j.
\]
The exact row reductions give $\dim K=15$ and $\dim N=16$.
Consequently the sum defining $N$ is direct. Here the direct sum is an
internal decomposition of modules on the same $120$ coordinates, not a
concatenation of codes on disjoint coordinate sets. That common lemma proves
self-orthogonality of both codes. These ranks and the full table, including
its underlying block lists, are obligations of the complete finite certificate.

\subsection{Recognizing the socle}
The literal permutations in \texttt{module\_data.json} preserve the residual Steiner graph and generate a primitive
permutation group $G\cong L_3(4)$ on $\mathcal H$, with point stabilizer
$L_3(2)$. The certificate checks primitivity, the isomorphisms and stabilizer, so this is
the specified degree-$120$ coset action in the source. It also checks
invariance of $K$ and $N$. We apply the common full-matrix-algebra and socle criteria \cite[Theorems~6.2--6.3]{U}; the whole filtration and all intertwiner spaces must be checked.

Here $V=\mathbb F_3^{120}$. The supplied basis certifies invariant prefixes
of dimensions
\[
0,15,34,49,64,65,84,99,100,119,120.
\]
The ten factor dimensions and homomorphism-space dimensions are
\[
\begin{array}{c|rrrrrrrrrr}
i&1&2&3&4&5&6&7&8&9&10\\\hline
\dim D_i&15&19&15&15&1&19&15&1&19&1\\
\dim\operatorname{Hom}_G(D_i,V)&1&0&0&0&1&0&1&1&0&1.
\end{array}
\]
Triangularity of the generator matrices verifies invariance of all prefixes.
For each factor the matrix algebra has dimension $(\dim D_i)^2$, proving
simplicity by the common criterion. The first prefix is exactly $K$.
For each factor, the verifier forms \emph{all} equations
$XP_g=D_gX$ for the given generators and computes their full solution space.
Every row of every basis solution belongs to $N$. Linearity then establishes
the hypothesis for every homomorphism, not only a selected map.
The constant line is simple and $K$ is the first simple factor, so
the common criterion yields
\[
 \Soc(\mathbb F_3^{120})=N=K\oplus\mathbb F_3j.
\]
This calculation does not classify every irreducible module of $L_3(4)$.
The source's classification assertion concerning three nonisomorphic faithful
irreducible modules of dimension $15$ is retained from its reference [13], Jansen--Lux--Parker--Wilson, \emph{An Atlas of Brauer Characters} (1995).
Our certificate identifies a particular such simple module; its faithfulness
also follows because $G$ is simple and the representation is nontrivial.

\subsection{Weights and full permutation automorphisms}
Complete enumeration in the certified bases gives
\[
\begin{array}{c|r|r|r|r|r}
\text{code}&\text{number of words}&\text{dimension}&d&\text{words of weight }d&d^\perp\\\hline
K&14348907&15&48&1260&4\\
N&43046721&16&48&1260&4.
\end{array}
\]
The enumerator visits each of the $3^k$ coefficient vectors exactly once;
independence of the basis makes these distinct words. It records the full
weight distribution. The complete dual distributions follow from the common MacWilliams theorem \cite[Theorem~7.1]{U}. 
In both
cases the dual coefficients at weights $1,2,3$ vanish. At weight $4$ they are
$2100$ for $K^\perp$ and $1260$ for $N^\perp$. Thus the dual dimensions are $105$
and $104$, respectively, and the stated dual minimum distances follow.
These are complete finite computations, not distance bounds inferred from
self-orthogonality.

The orbit under $G$ of the supplied weight-$48$ word contains $1260$ distinct
words and $630$ distinct supports. Since $K\subset N$ and both codes have
exactly $1260$ minimum words, this orbit is the complete minimum-word set of
both codes. Every coordinate-permutation automorphism of either code
therefore acts on its support-incidence structure, with $120$ coordinate
vertices and $630$ distinct support vertices.

The common refinement theorem \cite[Theorem~6.4]{U} bounds the full automorphism group of this structure. 
Successively individualize the coordinate vertices
$1,6,48,10,2,3$ in the literal certificate ordering. Their successive
refined colour classes have sizes
\[
120,21,2,4,2,2.
\]
After these choices all coordinate classes are singletons. The theorem gives the upper bound $120\cdot21\cdot2\cdot4\cdot2\cdot2=80640$.

The explicit permutations supplied in the certificate generate a group
$A$ of order $80640$ and preserve each of $K,N$. Thus both code automorphism
groups attain the upper bound and equal $A$. The certificate checks
$A'=G$ and supplies two commuting involutions generating a complement of
order $4$, disjoint from $G$. Consequently
\[
 \Aut(K)=\Aut(N)=A\cong L_3(4):2^2.
\]
No graph-isomorphism program is needed for this upper bound. Such software
was useful in discovering the permutations, but the retained proof checks
literal permutations, ordinary group orders and the full refinement chain.


\section{The irreducible code on 176 residual blocks}
Here $H=H_1$, $C=C_{176}=H^{\mathsf T}(e^\perp)$ and $\dim C=21$ by (1).
The integral Gram matrix is $40I_{22}+16J_{22}$, reducing to $I_{22}+J_{22}$.
The common incidence theorem proves LCD and the two complement--dual
identities. The constant vector belongs to $C^\perp$, because its image
under $H$ is $56e$, which annihilates $e^\perp$. Thus the full orthogonal
decomposition in the source is established before any module classification.

\subsection{The exact module and group certificate}
The literal permutations in \texttt{module\_data.json} preserve every
adjacency entry of $\Gamma$. Their closure, computed by exact permutation
composition, has order $443520$. They also permute the $22$ point-incidence
rows of $H$. The resulting faithful degree-$22$ action is explicitly
identified with \texttt{MathieuGroup(22)} by an isomorphism checked in GAP.
Thus the subgroup $G$ they generate is $M_{22}$. This conventional group
identification is distinct from proving that $G$ is the \emph{full} code
automorphism group. The latter will follow from the common refinement theorem \cite[Theorem~6.4]{U}.
The auxiliary group calculation also verifies that the degree-$176$ action
is primitive with stabilizer $A_7$ of order $2520$.

The row-reduced adjacency basis is checked byte for byte against the basis
of the retained complete enumeration. For each coordinate generator, its
images on these $21$ rows give a $21\times21$ representation matrix.
The checker reconstructs those images from the proposed matrices, so the
matrix module is the actual code $C$, not a representation with only the
same dimension. Deterministic multiplication and elimination produce
$441$ independent generator words, all of length at most eight. Each word
is evaluated again and the $441$ flattened matrices are row-reduced.
Their rank is $441$. A separate GAP calculation confirms the same generated
algebra dimension. The common matrix-algebra theorem \cite[Theorem~6.2]{U} proves that $C$ is absolutely
irreducible. This action is faithful as well: the $176$ columns of a code
basis are distinct, so a coordinate permutation fixing every codeword
fixes every coordinate. Column distinctness is checked in the certificate.

The retained complete enumeration gives $d(C)=56$ and exactly $924$ words
of weight $56$. The supplied literal weight-$56$ word belongs to $C$.
Its orbit under $G$ has $462$ words and is disjoint from the orbit of its
negative. Together those two orbits have $924$ words, so they exhaust the
minimum words. Their supports are the same $462$ distinct subsets.
These assertions are checked by finite orbits and the already accepted
count; no new exhaustive codeword computation is involved.

For this support-incidence graph, with $176+462$ vertices, the deterministic
refinement rule of the common refinement theorem \cite[Theorem~6.4]{U} successively individualizes
coordinates $1,5,29,61,2$. The corresponding cell sizes are
\[
 176,\quad70,\quad6,\quad2,\quad3.
\]
All coordinate colours then differ. Therefore
\[
 |\Aut(C)|\leq176\cdot70\cdot6\cdot2\cdot3=443520.
\]
Every graph automorphism preserves its adjacency row space, so
$G\leq\Aut(\Gamma)\leq\Aut(C)$. The already checked order of $G$ forces
equality throughout:
\[
 \Aut(C)=\Aut(\Gamma)=G\cong M_{22}.                 \tag{10}
\]
This upper bound uses no graph-isomorphism backend. Such software was
used only to discover the two literal generators.

The complete retained weight distribution determines the dual by the common MacWilliams theorem \cite[Theorem~7.1]{U}.
The certificate computes all coefficients exactly, checks nonnegative
integrality and total $3^{155}$, and obtains zero at weights one through
five and $184800$ at weight six. Thus the parameters are
$[176,21,56]_3$ and $[176,155,6]_3$, with the asserted minimum counts.

\subsection{The precise representation boundary}
Theorem~1(c) describes the \emph{unique $21$-dimensional} irreducible
$F_3$-module for $M_{22}$. Proposition~4's final sentence omits the dimension,
but its proof immediately restores it and invokes the known $3$-modular
character table. We retain that intended classification statement, and
the proof's further assertion that $21$ is the smallest nontrivial
irreducible dimension, from the source's reference [13]:
Jansen--Lux--Parker--Wilson, \emph{An Atlas of Brauer Characters} (1995).
These assertions are understood up to module isomorphism. Our full
matrix-algebra calculation proves irreducibility of this particular
module; it does not classify all modules. The unqualified literal
uniqueness sentence cannot be correct, since the trivial one-dimensional
module is irreducible too.

The source's historical near-optimality comparison for the dual is a
separate assertion based on its tables, not a current table claim or a
consequence of (10). No additional irreducibility is asserted for the
$155$-dimensional dual, which contains the constant line. The group in
(10) consists of coordinate permutations; allowing scalar monomial maps
would change the automorphism convention.


\section{The augmentation code on 253 heptads}

Let $\mathcal D$ be the specified $S(4,7,23)$, with point set $\Omega$
and $253$ heptads. Its point--block incidence matrix is $H$.
Two distinct heptads intersect in one or three points, and $\Gamma$ joins
them in the latter case. Write $A$ for its adjacency matrix. All vector
spaces in this note are over $F=\mathbb F_3$. Put
\[
 E=\{x\in F^{23}:\textstyle\sum_i x_i=0\},\quad
 C=\im A,\quad D=H^{\mathsf T}E.
\]
The shared construction checks the Steiner condition and actual incidence
matrix, not merely the strongly regular parameters.

\begin{theorem}\label{L050.253}
The code $C$ is the orthogonal direct sum $\langle\mathbf j\rangle\oplus D$,
where $D$ is an absolutely irreducible $22$-dimensional module for $M_{23}$.
Both $C$ and $D$ are LCD. Their parameters and minimum multiplicities are
\[
\begin{array}{c|c|r}
\text{code}&\text{parameters}&\text{minimum words}\\\hline
C &[253,23,77]_3&46\\
D &[253,22,97]_3&3542\\
C^\perp &[253,230,6]_3&850080\\
D^\perp &[253,231,6]_3&850080.
\end{array}
\]
The full groups of coordinate permutations preserving $C$ and $\Gamma$
are equal and isomorphic to $M_{23}$.
\end{theorem}

\begin{proof}
The common residual-incidence theorem \cite[Theorem~2.2]{U} applies to the fixed system and gives $HH^{\mathsf T}=2I_{23}$,
$A=H^{\mathsf T}(J_{23}-I_{23})H$, $C=\im H^{\mathsf T}$ and the
orthogonal splitting $C=Fj\perp\!\oplus H^{\mathsf T}(e^\perp)$.
The two point-space summands are nondegenerate because $23\ne0$ in $F$.
Thus it proves both LCD statements and the displayed dimensions.

We next prove the irreducibility, using explicit matrices in place of a
modular character-table inference. The five permutations of $\Omega$
in \texttt{generators.json} preserve the full heptad set. They therefore
act on the code coordinates and on $E$. In the basis
$e_i-e_{23}$, $1\leq i\leq22$, the row matrix for a point permutation
$g$ has entries
\[
 M(g)_{ij}=\mathbf1_{j=g(i)}-\mathbf1_{j=g(23)},\qquad 1\leq i,j\leq22.
\]
An indicator with point $23$ is zero because $j\leq22$.
The certificate checks that multiplication by these matrices agrees
with the induced permutation of all $253$ coordinates on the actual
rows $h_i-h_{23}$. It produces $484$ independent words in these matrices,
each of length at most six, and reevaluates every word before reducing
the flattened matrices. They span $M_{22}(F)$.

The common full-matrix-algebra criterion \cite[Theorem~6.2]{U} proves absolute irreducibility of $E$ and its equivariant image $D$.

The existing enumeration records contain the complete weights of the
row-reduced adjacency code and of the difference code. Their bases are
compared byte for byte with the bases reconstructed here. The latter
row space is $D$: its dimension is $22$, it lies in $C$ and is orthogonal
to $\mathbf j$. The recorded enumerations have respectively $3^{23}$
and $3^{22}$ coefficient vectors. They give the two distances and counts
in the first two rows of the table. The dual distributions follow from the common MacWilliams theorem \cite[Theorem~7.1]{U}.
All coefficients are evaluated over the integers. Nonnegativity,
integrality and total $3^{253-k}$ are checked. Both duals have no words
of weights one through five and exactly $850080$ of weight six.
This proves the table without repeating either exhaustive enumeration.

The $46$ words $h_i,-h_i$ are precisely the minimum words, by their weights and the complete count. The common Steiner recognition proposition \cite[Proposition~2.3]{U} therefore identifies $\Aut(C)$ with the block action of $\Aut(\mathcal D)$. This explicitly discharges its minimum-word hypothesis.

For an explicit upper bound, form the bipartite point--block incidence
graph, with its two parts initially given different colours. Refine the
colours by the old colour and the counts of neighbours of each colour.
After individualizing points $1,2,3,4,5,6$ successively, the certificate
obtains current cell sizes
\[
 23,\quad22,\quad21,\quad20,\quad16,\quad3.
\]
All points are then distinguished. A pointwise automorphism fixes every
heptad, since distinct heptads have different point neighbourhoods.
The common refinement theorem \cite[Theorem~6.4]{U} consequently bounds $|\Aut(\mathcal D)|$ by $23\cdot22\cdot21\cdot20\cdot16\cdot3=10200960$.
Exact group computations on the literal generators give a subgroup of
this order and an isomorphism to the standard group $M_{23}$.
The induced block action is checked to be faithful. Thus the upper bound
is attained and identifies the full design and code groups. Finally the
subgroup preserves $\Gamma$, and every graph automorphism preserves its
adjacency row space. Therefore
$G\leq\Aut(\Gamma)\leq\Aut(C)=G$, proving the asserted equality.
The upper bound uses only incidence refinement; graph-isomorphism
software was used to discover the generators and is not a final premise.
\end{proof}

\paragraph{Qualifications to Proposition 5.}
The main Theorem~1(d) is recovered as stated. Proposition~5 additionally
assigns distance eight to $D^\perp$ and declares it irreducible. Both
assertions are false for these codes. Since $D\subset C$, the invariant
space $C^\perp$ of dimension $230$ is a proper nonzero subspace of
$D^\perp$, whose dimension is $231$. The same inclusion already gives
$d(D^\perp)\leq d(C^\perp)=6$; the exact value six is verified above.
The proposition also calls $D$ the ``unique and smallest'' irreducible
module without qualifications. The unqualified statement overlooks the
trivial one-dimensional module. Its intended classification requires an
input separate from our matrix certificate. We retain the complete
$3$-modular character table of $M_{23}$, as in the source's reference [13],
Jansen--Lux--Parker--Wilson, \emph{An Atlas of Brauer Characters} (1995).
Its irreducible degrees over an algebraic closure are
\[
 1,22,45,45,104,104,231,253,770,770,990,990,1035.
\]
The accompanying GAP calculation checks these entries against its
installed table \texttt{M23mod3}; it does not reprove completeness of
that table. Together with standard scalar extension and descent for
simple modules over finite fields, the table gives a unique
$22$-dimensional irreducible $F$-module, and $22$ is the least
nontrivial degree. Indeed a simple module splits after scalar extension
into an orbit of absolutely simple constituents of equal degree. Below
degree $22$ the only available constituent is trivial, which forces the
original simple module to be trivial. At degree $22$ the unique
constituent is the one already realized by $D$; descent of module
isomorphisms gives uniqueness over $F$. These are retained
classification premises, explicitly distinct from the direct proof of
irreducibility of $D$. The source's open uniqueness question for the
strongly regular parameters $(253,140,87,65)$ is not resolved here.


\section{The source questions and precise qualifications}

Section~6 asks whether equality of a graph code with the orthogonal dual
of its complementary graph code forces LCD, and whether the converse
holds. The forward direction holds over every field: that equality gives
$A(J-I-A)=0$, and the common projection theorem
\cite[Theorem~1.2]{U} gives LCD. Its conclusion also identifies the exact
one-dimensional defect when both degrees vanish in the field.

The converse fails already for the four-cycle over $\mathbb F_3$.
In cyclic coordinate order its code is spanned by
$(1,0,1,0)$ and $(0,1,0,1)$. Their Gram matrix is $2I_2$, so this code is
LCD. Its complementary graph is two disjoint edges, whose adjacency
matrix is invertible. The complementary code is the whole space, whose
dual is zero, and hence cannot equal the two-dimensional cycle code.
The examples of lengths $120$ and $176$ still give the affirmative
existence answer sought in the source's introduction: they go beyond
triangular graphs and have exact complementary dual codes. These
arguments make no claim about the historical priority of either implication.

The uniqueness problem for strongly regular parameters $(253,140,87,65)$,
mentioned in Section~6.3, remains outside our result. Our $253$-coordinate
code is constructed from the specified Steiner system, and no uniqueness
of every graph with those parameters is used. Remark~3's discussion of
locating the other two $15$-dimensional irreducible modules is also not a
construction of those modules here. Their abstract existence and
classification remain the explicitly retained Brauer-table input; the
specified degree-$120$ permutation module and its particular socle are
proved above.

The asymptotic Gilbert--Varshamov result of Sendrier~\cite{Sendrier}
concerns families of LCD codes. Membership of one finite code in the LCD
class does not by itself assert a finite optimality comparison or an
asymptotic property of that single code. We retain the cited asymptotic
result with that meaning. Likewise, the comparisons to code tables in
Propositions~1--4 and Corollary~1 are historical comparisons, not updated
optimality claims. None is used to establish a dimension, distance or
full automorphism group here.

The literal and intended readings of the representation statements have
been separated throughout. The $105$ global least-degree statement is
false without an additional restriction, as witnessed by the actual
nontrivial one-dimensional quotient character. The dimension qualifier
in the $176$ main theorem is necessary. For length $253$ the printed
$231$-dimensional dual has distance six rather than eight and is reducible;
its invariant $230$-dimensional subcode already disproves irreducibility.
The main Theorem~1(d), which does not assert those erroneous dual claims,
is preserved. A matrix-algebra certificate proves irreducibility of a
specified representation; global uniqueness and minimum-degree
classifications require the separately retained table inputs.

\paragraph{Computations and formal scope.}
The finite premises are the complete specified constructions, literal
matrices and group generators, full filtration and intertwiner systems,
complete weight certificates, and support-refinement upper bounds.
Their executable files and mathematical data are inventoried in full.
Discovery with graph-isomorphism software or MeatAxe is not a replacement
for these final checks. Exact arithmetic in the declared Python, C and
GAP implementations is trusted in the finite certificates. Operational
drivers and optional weakening tests are listed separately. The nine
previously checked generic Lean declarations for projection and averaging
retain their precise restricted scope; this application is not an
end-to-end Lean formalization. No new source, question, historical novelty,
or compression credit is assigned by this author draft.

\appendix
\section{The two complete distributions on 105 coordinates}
This table reproduces both nonzero supports of the weight distributions
in source Tables~1--2, using the exact accepted enumeration. All omitted
weights have multiplicity zero. The zero word is included. The totals are
$3^{20}$ and $3^{19}$, respectively.
\begin{center}
\begin{tabular}{r|r|r}
\toprule
Weight&$C_{105}$&$D_{105}$\\\midrule
0&1&1\\
33&770&0\\
36&3360&3360\\
42&29640&9480\\
45&38304&336\\
48&141120&100800\\
51&987840&144480\\
54&3649800&1659840\\
57&21993720&6552000\\
60&108522288&36809136\\
63&303515520&101499720\\
66&606412800&202063680\\
69&818230000&269951640\\
72&804875190&272891220\\
75&526684032&172986912\\
78&222579280&74177040\\
81&58508240&20054720\\
84&9106500&2743020\\
87&1230880&544320\\
90&127008&54432\\
93&128520&0\\
96&15680&15120\\
99&3360&0\\
105&548&210\\
\bottomrule
\end{tabular}
\end{center}


\begin{thebibliography}{9}
\bibitem{LR} D.~Leemans and B.~G.~Rodrigues,
Linear codes with complementary duals from some strongly regular
subgraphs of the McLaughlin graph, \emph{Math. Commun.}~21 (2016), 239--249.
\bibitem{U} \emph{Incidence operators, orthogonal codes and permutation modules},
fixed common manuscript, additive edition~2; theorem labels and exact
edition are bound in the accompanying application manifest.
\bibitem{Coolsaet} K.~Coolsaet, The uniqueness of the strongly regular
graph srg$(105,32,4,12)$, \emph{Bull. Belg. Math. Soc. Simon Stevin}~12
(2005), 707--718, Theorem~3; retained as in \cite[reference~6]{LR}.
\bibitem{Atlas} C.~Jansen, K.~Lux, R.~Parker and R.~Wilson,
\emph{An Atlas of Brauer Characters}, Clarendon Press, Oxford, 1995;
retained classification inputs as in \cite[reference~13]{LR}.
\bibitem{Sendrier} N.~Sendrier, Linear codes with complementary duals
meet the Gilbert--Varshamov bound, \emph{Discrete Math.}~285 (2004),
345--347; retained asymptotic statement.
\end{thebibliography}
\end{document}
